\documentclass[10pt,conference]{IEEEtran}
\usepackage{cite}
\usepackage{amsmath,amssymb,amsfonts}
\usepackage{algorithmic}
\usepackage{graphicx}
\usepackage{textcomp}
\usepackage{xcolor}

\begin{document}

\title{Real-Time Driver Drowsiness Detection via Adaptive Geometry and Signal Reliability Gating}

\author{\IEEEauthorblockN{Sayemuddin}
\IEEEauthorblockA{\textit{Department of Computer Science \& Engineering} \\
\textit{Lead AI Research Engineer}\\
Email: author@domain.com}
}

\maketitle

\begin{abstract}
Driver fatigue and drowsiness are leading contributors to global vehicular accidents. While facial landmark heuristics (e.g., Eye Aspect Ratio) provide high-speed non-invasive drowsiness detection, they suffer from false positives due to lighting fluctuations, camera shake, and individual biometric variations. Conversely, deep convolutional neural networks offer high classification accuracy but incur unacceptable computational overhead on resource-constrained embedded automotive hardware. In this paper, we present a lightweight, real-time driver drowsiness detection architecture whose central component is a \emph{decomposed signal-reliability gate} that multiplicatively attenuates multi-cue fatigue evidence by a continuous, component-wise signal-quality index before temporal accumulation; a separate state-level guard governs exit from the severe-fatigue state. A complementary speech-jitter filter, gating on the mean absolute per-frame change in the mouth aspect ratio, suppresses mouth-motion false positives. We evaluate on NTHU-DDD under leave-one-subject-out cross-validation, measuring false-positive rate at a matched detection rate. Under this strict subject-disjoint protocol we report an \emph{honest negative result} at the frame level: the reliability gate does not reduce the false-positive rate at a matched true-positive rate of $0.80$, and the speech-jitter filter increases it. We analyze why a frame-level operating-point metric on clip-derived labels cannot observe the gate's episode-level, spurious-alarm–suppression behaviour, and we establish real-time host feasibility ($3.205$~ms/frame on a laptop-class ARM CPU); on-device Raspberry~Pi~4 profiling remains future work.
% RESULTS POLICY: every quantitative figure in this manuscript must trace to a
% logged run in EXPERIMENT_REGISTRY.md and a committed artifact under results/
% or experiments/. No Raspberry Pi 4 number exists yet; do not insert one.
\end{abstract}

\begin{IEEEkeywords}
Driver Drowsiness Detection, Signal Reliability Gating, Eye Aspect Ratio, Embedded Vision, TensorFlow Lite, Edge Computing.
\end{IEEEkeywords}

\section{Introduction}
Driver drowsiness detection systems (DDDS) are vital for intelligent transportation safety. Classical geometric approaches monitor facial landmark aspect ratios such as EAR and MAR. However, fixed thresholds fail across subject demographics and suffer under severe illumination degradation.

To address these challenges, we propose:
\begin{enumerate}
    \item A \textbf{decomposed signal-reliability gate} that combines per-frame landmark-stability, brightness, and cue-consistency sub-scores into a continuous reliability index, and multiplicatively attenuates the fused fatigue evidence \emph{before} temporal accumulation for all states; a separate state-level guard governs exit from the severe-fatigue state.
    \item A \textbf{speech-jitter filter} that suppresses mouth-motion (talking) false positives by gating on the mean absolute per-frame change in the mouth aspect ratio.
    \item A \textbf{2D image-plane geometry formulation} that computes both EAR and MAR from image-plane landmark coordinates to avoid monocular depth inflation, reported as a supporting engineering result.
\end{enumerate}
An optional selective-CNN eye-validation arm (MicroEyeNet) is evaluated only as an ablation and is not claimed as a contribution, as selective CNN invocation is prior art.

\section{Proposed Architecture}
Our architecture consists of the following stages:
\begin{itemize}
    \item \textbf{Facial Landmark \& Geometry Processing}: MediaPipe FaceMesh landmarks; EAR and MAR both computed from image-plane landmark coordinates (2D) to avoid monocular depth inflation.
    \item \textbf{Head Pose Estimation}: pitch/yaw/roll via \texttt{solvePnP}.
    \item \textbf{Signal-Reliability Estimation}: per-frame landmark-stability, brightness, and cue-consistency sub-scores combined into a continuous reliability index.
    \item \textbf{Reliability-Gated Fusion}: multiplicative attenuation of fused fatigue evidence before temporal accumulation, applied to all states; a separate state-level guard governs exit from the severe-fatigue state.
    \item \textbf{Asymmetric Temporal Integration \& Hysteresis State Machine}: wall-clock (monotonic) time-based integration, FPS-independent.
    \item \textbf{Speech-Jitter Filter}: gating on the mean absolute per-frame change in the mouth aspect ratio to suppress talking false positives.
\end{itemize}

\section{Experimental Results}
% All figures below are backed by logged experiments (EXP-001..004) in
% EXPERIMENT_REGISTRY.md and committed artifacts under experiments/ and
% results/. Do not add numbers that lack an EXP-### row + artifact.
All numbers reported here are drawn from logged experiments
(\mbox{EXP-001}--\mbox{EXP-004}) with committed artifacts; no unmeasured value
is cited.

\subsection{Eye-state model (MicroEyeNet) and quantization}
The optional MicroEyeNet eye-state validator ($19{,}745$ parameters) was trained
on a genuinely subject-disjoint split of the MRL Eye dataset (\mbox{EXP-002};
seed~42), reaching a test $F_1$ of $0.9623$. Post-training INT8 quantization
(\mbox{EXP-003}) compresses the model to $25.55$~KB ($3.18\times$ smaller) with a
$-0.026\%$ change in $F_1$, yielding a deployment-ready asset. The CNN arm is
evaluated only as an ablation and is not claimed as a contribution.

\subsection{Ablation on NTHU-DDD (leave-one-subject-out)}
We evaluate five frozen variants under subject-disjoint LOSO on the full
NTHU-DDD corpus ($66{,}521$ frames, $4$ subjects; \mbox{EXP-004}), toggling the
speech-jitter filter, the reliability gate, and the CNN arm: V0 (baseline
weighted fusion), V1 (\,$+$\,speech filter), V2 (\,$+$\,reliability gate), V3
(full), and V4 (full\,$+$\,CNN). The operating point is fixed on the V0 ROC at a
matched true-positive rate of $0.80$ and held constant across variants; the
sweep variable is the continuous per-frame fatigue score.

\begin{table}[t]
\caption{Frame-level LOSO ablation on NTHU-DDD (EXP-004). FPR is reported at a
matched TPR of $0.80$ fixed on the V0 baseline. Lower FPR and higher AUC are
better.}
\label{tab:ablation}
\centering
\begin{tabular}{lccc}
\hline
Variant & Gate / Filter / CNN & ROC-AUC & FPR@TPR$=0.80$ \\
\hline
V0 & --\,/\,--\,/\,-- & $0.6288$ & $0.6241$ \\
V1 & --\,/\,\checkmark\,/\,-- & $0.6170$ & $0.6694$ \\
V2 & \checkmark\,/\,--\,/\,-- & $0.6245$ & $0.6244$ \\
V3 & \checkmark\,/\,\checkmark\,/\,-- & $0.6133$ & $0.6690$ \\
V4 & \checkmark\,/\,\checkmark\,/\,\checkmark & $0.6133$ & $0.6690$ \\
\hline
\end{tabular}
\end{table}

Table~\ref{tab:ablation} reports an \textbf{honest negative result}. At the
matched operating point, the reliability gate in isolation (V2) leaves the
false-positive rate essentially unchanged relative to the baseline ($0.6244$ vs.
$0.6241$), and every variant that adds the speech-jitter filter (V1, V3, V4)
\emph{raises} the false-positive rate to $\approx 0.669$. ROC-AUC never exceeds
the baseline (V0 highest at $0.6288$). The proposed mechanisms therefore do not
improve frame-level discrimination under this protocol.

Two structural observations explain the result and constrain its
interpretation. First, V4 is byte-identical to V3: the CNN verdict is consumed
only by the alarm-decision stage and never modifies the swept fatigue score, so
it cannot move a frame-level ROC by construction. Second, NTHU-DDD labels are
clip-condition derived, making the frame-level false-positive rate a
conservative estimator: an open-eyed frame within a ``yawning'' clip is still
labelled drowsy. The gate's intended benefit --- suppression of \emph{spurious
alarm episodes} under unreliable measurement conditions --- and the CNN's
alarm-suppression effect are episode-level phenomena that a frame-level
operating-point metric cannot observe. Motivated by this analysis, we carried
out an event-level alarm evaluation of the same five variants on NTHU-DDD: it
confirms the negative at the event level, with an event recall of $0.122$ at
$6.5$--$9.7$ false alarms per hour and alarms firing for only $2$ of the $4$
subjects, so the gate does not deliver the episode-level spurious-alarm
suppression the design targets.

\subsection{Real-time feasibility}
End-to-end per-frame latency was measured at $3.205$~ms (median $3.080$~ms, 95th
percentile $3.316$~ms; $\approx 312$~FPS) over NTHU-DDD frames on a laptop-class
ARM CPU (\mbox{EXP-001}). This establishes host feasibility only; \textbf{no
Raspberry~Pi~4 measurement exists}, and no on-device latency is claimed.

\section{Conclusion}
We presented a signal-reliability-gated driver drowsiness detection
architecture optimized for embedded deployment, and evaluated it under strict
subject-disjoint leave-one-subject-out cross-validation. Our central empirical
finding is a negative one: at the frame level, on clip-derived NTHU-DDD labels,
the reliability gate does not reduce the false-positive rate at a matched
detection rate, and the speech-jitter filter increases it. We showed this
outcome is partly structural --- the CNN arm cannot influence the frame-level
score, and clip-derived labels make frame-level false positives a conservative
estimator --- and argued that the gate's intended benefit is an episode-level
effect --- a hypothesis our event-level evaluation did not bear out under the
current design. Future work is therefore on-device Raspberry~Pi~4 profiling and
on-road validation.
% Internal traceability (not for publication): the future work above maps to the
% official experiment roadmap in EXPERIMENT_REGISTRY.md \S4 --- event-level
% evaluation = EXP-005, gate redesign = EXP-006, Raspberry Pi deployment =
% EXP-007, optional second-dataset validation = EXP-008. EXP-### are internal
% registry IDs and deliberately do not appear in the manuscript body.

\bibliographystyle{IEEEtran}
\end{document}
